\appendix

\section{Proof of Theorem~\ref{thm:bridge}}\label{app:bridge}

Throughout, labels are $a=(\theta,i,t)\in\calS$, and $k_a,\beta_a,\calF_a,\J,\phi_U,\Gam,\rhoHG$ and program \eqref{eq:lp} are as in Sections~\ref{sec:setting}--\ref{sec:certificates}.

\paragraph{Step 1 (control-moment inequality).}
For a block $J_j$ and any measurable $u$ with values in $U$,
\[
\Bigl(\int_{J_j}k_a\Bigr)^{\!\top}\frac1{|J_j|}\int_{J_j}u-\int_{J_j}k_a^\top u=\int_{J_j}(\bar k_{a,\calH}-k_a)^\top u\le\int_{J_j}h_U(\bar k_{a,\calH}-k_a),
\]
using $\sup_{v\in U}w^\top v=h_U(w)$. Summing over blocks proves \eqref{eq:outer}; the averages remain compatible with the information structure because the mesh contains all sensing times, so common-history controls use the same block variable. For the uniform bound, choose a centre $u_c$ of $U$: since $\int_{J_j}(\bar k-k)=0$ the centre cancels and $e_{a,\calH}\le R_U\|k_a-\bar k_{a,\calH}\|_{L^1}$; a function of bounded variation satisfies $\|k_a-\bar k_{a,\calH}\|_{L^1}\le h_u\TV(k_a)$ (for every $\sigma\in J$, $\|\bar k-k(\sigma)\|\le\TV(k;J)$, the norm of any increment being at most the total variation; integrate over $J$ and sum), whence $e_{a,\calH}\le R_UV_Uh_u$ under (H2), (H4).

\paragraph{Step 2 (lower bound of \eqref{eq:sandwich}).}
Take any admissible policy and form its block averages, which lie in $U$ and satisfy \eqref{eq:outer}. For every grid label, $A_av-\beta^+_a-\ebar_a\le\int_0^Tk_a^\top u^\pi_\theta-\beta_a=\calF_a(\pi)$, using only the one-sided bound of \eqref{eq:outer} and $\beta^+_a\ge\beta_a$ (the reflected error can be larger than $\ebar_a$, so no absolute-value bound is inferred). The block-moment vector of every admissible policy is thus feasible for \eqref{eq:lp} at objective $\max_a\calF_a(\pi)$; minimising over policies gives $\rhoHG\le\J$.

\paragraph{Step 3 (upper bound of \eqref{eq:sandwich}).}
Take an optimiser $v^\ast$ of \eqref{eq:lp} and implement the corresponding piecewise-constant information-adapted control, for which $\int_0^Tk_a^\top u^{v^\ast}_\theta=A_av^\ast$ exactly: the implemented control is constant on each block, so no moment error arises; the price is paid only through grid labels versus all times. At grid times $\calF_a(v^\ast)\le\rhoHG+\ebar_a+\delta_\beta$, and (H1) extends the bound to every time within $h_t$ of a grid label, giving $\max_a\calF_a(v^\ast)\le\rhoHG+\max_{\hat a\in\calS_G}\ebar_{\hat a}+\delta_\beta+Lh_t$. Since $\J$ minimises over all policies, the upper bound follows; the uniform choice $\ebar\equiv R_UV_Uh_u$ (Step 1) gives the special case, and \eqref{eq:distbound} the sharpened budget. The $\max_a\ebar_a$ form is forced: the held control's exit from the continuous value is charged by the error bounds the program actually used.

\paragraph{Step 4 (sound finite certificates).}
Write the block input constraints as $F_{\mathrm{blk}}v\le f_{\mathrm{blk}}$. The dual of \eqref{eq:lp} is
\begin{equation}\label{eq:dual}
\max_{\lambda,\mu}\ -\lambda^\top(\beta^++\ebar)-\mu^\top f_{\mathrm{blk}}\quad\text{s.t.}\quad A^\top\lambda+F_{\mathrm{blk}}^\top\mu=0,\ \ \lambda\ge0,\ \mathbf 1^\top\lambda=1,\ \mu\ge0 ,
\end{equation}
so a positive dual objective certifies infeasibility of $\{Av\le\beta^++\ebar,\ v\in U^Q\}$. By Step 1 every genuinely safe measurable controller would generate a feasible $v$; the finite infeasibility certificate therefore excludes every such controller. Moreover the weighted continuous-time violation is at least the dual objective: a normalised certificate with margin $\gamma>0$ guarantees an actual facet violation of at least $\gamma$ in one stored test (Lemma~\ref{lem:semantics}).

\paragraph{Step 5 (completeness).}
If the system is nonviable, let $\gamma=\J>0$ and choose meshes and stored scenarios making the inflation in \eqref{eq:sandwich}---including the stored-scenario widths $\delta_{\mathrm{init}},\delta_{\mathrm{num}}$, which near-maximising stores contribute and completeness must drive to zero---smaller than $\gamma$; then $\rhoHG>0$ and Step 4 yields a positive certificate. With exactly stored scenarios and common width $h$, the inflation is at most $(R_UV_U+R_DV_D+L)h$, so $\rhoHG\ge\gamma/2$ as soon as that quantity is below $\gamma/2$; enumerating admissible meshes and halting on the first positive certificate is the semidecision procedure. No belief-state recursion is used.

\paragraph{Step 6 (exact atomic dual representation).}
\emph{(a) Policy compactness.} Each factor $\{u\in L^\infty(I_r;U)\}$ is weak-$\ast$ closed---$U$ is a finite intersection of half-spaces $f_i^\top v\le g_i$ and $\{u:f_i^\top u(\sigma)\le g_i\text{ a.e.}\}$ is weak-$\ast$ closed---and bounded; Banach--Alaoglu makes it weak-$\ast$ compact and convex, and the finite product $\Pol$ is compact, convex and metrisable.
\emph{(b) Joint continuity.} Each $\calF_a$ is a finite sum of $L^1$-pairings minus a constant, so $\pi\mapsto\calF_a(\pi)$ is weak-$\ast$ continuous. Under (H3), $\sup_\pi|\calF_a(\pi)-\calF_{a'}(\pi)|\le M_U\|k_a-k_{a'}\|_{L^1}+|\beta_a-\beta_{a'}|\to0$ with $M_U=\sup_{v\in U}\|v\|$ (H\"older), uniformly in $\pi$; with separate continuity this gives joint continuity of $(a,\pi)\mapsto\calF_a(\pi)$ on the compact $\calS\times\Pol$. The $L^1$ form of (H3) is the operative one: sup-norm continuity is false for the step kernels, which jump at $\sigma=t$ for every $t$.
\emph{(c) Mixed extension and Sion.} Let $\calP(\calS)$ be the Borel probability measures on $\calS$ with the weak-$\ast$ topology, compact convex and metrisable by Riesz. The bilinear $G(\pi,\mu)=\int\calF_a(\pi)\dd\mu(a)$ is affine in each variable---hence quasi-convex and quasi-concave---and continuous in each (dominated convergence in $\pi$; pairing against the continuous $a\mapsto\calF_a(\pi)$ in $\mu$). Sion's theorem \cite{Sion1958} gives $\min_\pi\max_\mu G=\max_\mu\min_\pi G$ with the outer maximum attained, and $\max_\mu G(\pi,\cdot)=\max_a\calF_a(\pi)$ because the integral of a continuous function against a probability measure lies in its range while Dirac measures attain the ends. Hence $\J=\max_\mu\min_\pi G$.
\emph{(d) Continuity of the value function.} $g(\mu):=\min_\pi G(\pi,\mu)$ is weak-$\ast$ continuous: by (b) the family $\{a\mapsto\calF_a(\pi)\}_\pi$ is uniformly bounded and equicontinuous in $C(\calS)$, so Arzel\`a--Ascoli makes weak-$\ast$ convergence $\mu_n\to\mu$ force $\sup_\pi|G(\pi,\mu_n)-G(\pi,\mu)|\to0$.
\emph{(e) Atomic approximation.} Finitely supported measures are weak-$\ast$ dense in $\calP(\calS)$ (quantise $\calS$ into Borel sets of diameter below $1/n$ and move each set's mass to a representative point); for $\mu_n\to\mu^\ast$, the maximiser of (c), (d) gives $g(\mu_n)\to\J$.
\emph{(f) The $\Gamma$ representation.} For $\mu=\sum_\ell\lambda_\ell\delta_{a_\ell}$ the inner minimum decomposes over information cells $(r,C)$: it is the infimum over $u_{r,C}\in L^\infty(I_r;U)$ of $\int_{I_r}g_{r,C}^\top u_{r,C}$. The inequality $\ge\int\phi_U(g_{r,C})$ is pointwise; the converse is measurable selection---for a.e.\ $\sigma$ the minimisers of $v\mapsto g_{r,C}(\sigma)^\top v$ over the polytope form a face, and the lexicographically first minimising vertex in a fixed order is measurable in $\sigma$---so the infimum equals $\int_{I_r}\phi_U(g_{r,C})$, and summing over cells gives $\min_\pi\sum_\ell\lambda_\ell\calF_{a_\ell}(\pi)=\Gam(\lambda,a)$, which also proves Lemma~\ref{lem:semantics}.
\emph{(g) Nonviability equivalence.} If $\J>0$, (e) supplies a finitely supported $\mu$ with $g(\mu)>\J/2>0$; conversely $\Gam(\lambda,a)\le\min_\pi\max_a\calF_a(\pi)=\J$ because the maximum bounds the weighted average. This proves \eqref{eq:atomic} and Corollary~\ref{cor:dual}.

\paragraph{Step 7 (sparsity).}
Factor the finite row map through its rank, $Av=L(Rv)$ with $Rv\in\R^r$; the attainable moment set $\calM=\{Rv:v\in U^Q\}$ is compact and convex and each safety row defines a half-space in $\R^r$. If the common feasible set is empty, Helly's theorem in $\R^r$ gives an obstruction involving at most $r+1$ safety half-spaces, and finite LP duality on that subproblem \cite{Chvatal1983} yields a sparse positive certificate with the same objective: $r+1\le mQ+1$ labels. To preserve a prescribed margin $\gamma_0<\rhoHG$, apply the argument to the still-infeasible shifted system $A_av\le\beta^+_a+\ebar_a+\gamma_0$. The continuous finite-rank case follows identically on the compact attainable moment body of the lifted kernels, giving $r_\calI+1$. \qed

\section{Proofs for Section~\ref{sec:ladder}}\label{app:ladder}

\begin{proof}[Proof of Proposition~\ref{prop:ladder}]
(i) is the pooling identity: $\sum_j\lambda_j(n_j^\top w)=(\sum_j\lambda_jn_j)^\top w=0$ cannot have every summand strictly negative. (ii) The objective $\max(n_i^\top w,n_j^\top w)$ is piecewise linear, so its minimum over the hexagon is attained at a vertex or on the equal-projection segment $(n_i-n_j)^\top w=0$. On that segment both authorities coincide: $w=-c(1,2)$ for $\{1,2\}$, where the facet $|u_2|\le\tfr45$ binds at $c=\tfr25$, and its mirror for $\{1,3\}$; $w=(c,0)$ for $\{2,3\}$, where the vertex $n_1$ binds at $c=1$ and gives $\tfr35$. The vertices themselves give authority at most $\tfr7{25}$ for the mixed pairs (at $w=n_3$, resp.\ $w=n_2$) and $\tfr35$ for $\{2,3\}$ (at $w=n_1$). For the triple, (i) gives $0$, attained at $w=0$. (iii) Along the schedule the window exit state is $\tfr{34}{25}+\tau-t^\ast\tau^2/2$ at critical velocity $1-t^\ast\tau$, and full braking adds $(1-t^\ast\tau)^2/2$; collecting terms gives the displayed quadratic, whose vertex in $\tau$ sits at $1/t^\ast\ge1$, so it increases on the relevant range. Setting it equal to $2$ gives $6\tau^2-30\tau+7=0$ for $t^\ast=\tfr25$ and $6\tau^2-20\tau+7=0$ for $t^\ast=\tfr35$, whose smaller roots are the displayed thresholds in $\mathbb Q(\sqrt{183})$ and $\mathbb Q(\sqrt{58})$. The singleton peak $\tfr{93}{50}$ is delay-independent.
\end{proof}

\begin{proof}[Proof of Theorem~\ref{thm:pairs}]
Write the pair as $B=\{i,j\}$ and let $t=\tfr25$ for a pair containing branch~1 and $t=\tfr35$ for $\{2,3\}$. Give the two critical position facets the weights $(\tfr38,\tfr58)$ for $(1,2)$ or $(1,3)$ and $(\tfr12,\tfr12)$ for $(2,3)$. The pooled vector $\nu=\lambda_in_i+\lambda_jn_j$ is $(0,\tfr12)$, $(0,-\tfr12)$ or $(-\tfr35,0)$ respectively, and its support value on the hexagon is $h_U(-\nu)=t$ (direct evaluation at the six vertices); each individual normal has support value one. Choose both labels at the common time $s_\ast=\tau+q$, $q=1-t\tau>0$, with margin $\beta=\tfr{16}{25}-s_\ast$. The exact dual is \eqref{eq:pairdual}; it minimises over \emph{every} measurable common blind control and every independently chosen post-revelation control (Lemma~\ref{lem:semantics}), so $\Gamma_B>0$ is an all-control nonviability certificate. $\Gamma_B$ is strictly increasing for $0\le\tau<1/t$ and its first zero is the displayed root; at each root $s_\ast=1+(1-t)\tau<\tfr65$ (the rational brackets $0.2453<\tau_{12}<0.2454$ and $0.3973<\tau_{23}<0.3974$ verify this), so for every delay just above a root, still with $s_\ast\le T$, the dual is strictly positive. For an arbitrary delay beyond the root, choose an intermediate delay with positive dual and $s_\ast\le T$; its obstruction transfers to the later delay by Lemma~\ref{lem:mono}.

For the converse, at the root use the blind controls of Proposition~\ref{prop:ladder}: $w=(-\tfr25,-\tfr45)$ (or its reflection) or $w=(1,0)$, followed on branch $j$ by $u=-n_j$ for $q=1-t\tau$ and then $u=0$ until $T$. The critical position is nondecreasing until the projected velocity vanishes at $s_\ast$, where it equals $2$ by the quadratic identity, and then remains constant. For every other position facet $i\ne j$, $n_i^\top n_j<0$ and $h_U(n_i)=1$, so for $0\le s\le\tfr65$, $n_i^\top p_j(s)\le(\tfr{34}{25}+s)\,n_i^\top n_j+s^2/2\le\tfr{18}{25}<2$. Each velocity coordinate is affine on each of the three control legs, so its largest absolute value occurs at a leg endpoint: initially $n_j$, at the blind endpoint $n_j+\tau w$, at the end of braking $\tau(w+tn_j)$; substituting the three $w$ and the rational upper root bounds shows every coordinate at each junction has magnitude at most $1$ (at most $\tfr45$ for $\{2,3\}$), strictly below $\tfr65$. These bounds verify all seven facets continuously, including the final hold leg (SM~S3 lists the junction values). An earlier revelation may be ignored until the certified root, so the schedule is admissible and safe for every smaller delay as well. This establishes the closed-endpoint thresholds, not merely isolated rational viability samples.
\end{proof}
